\documentclass{article}
\usepackage{graphicx}
\usepackage{float}
\usepackage{amsmath}
\usepackage{caption}
\usepackage{tabularx}
\usepackage{booktabs}
\usepackage[en-GB]{datetime2}
\DTMlangsetup[en-GB]{ord=omit}
\graphicspath{{./images}}
\captionsetup{textfont=small, labelfont=bf}
\renewcommand{\thesubsection}{\arabic{subsection}}


\title{University of Aberdeen\\Designing Real World Trials - PU5068\\Assessment 1 - Flow Diagram}
\author{Fernaldy Aristo Wirjowerdojo\\52642998}
\date{2 November 2026}

\begin{document}
\maketitle

\pagebreak
\section*{Rationale}
The study wants to compare the effectiveness of long-term GnRH analogues with add-back HRT against laparoscopic surgery in terms of clinical and cost for managing recurrent pain in women with endometriosis.
\subsection{PICOT}
The PICOT of this study is as follows:
\begin{itemize}
    \item P: Women aged between 21-49 years with pain recurrence following previous surgical treatment of endometriosis but wish to avoid removal of ovaries and hysterectomy
    \item I: Long term use of GnRH analogues with add-back HRT
    \item C: Laparoscopic ablation/excision of endometriosis
    \item O: Better functional health and well-being measured by EHP-30 and EQ-5D
    \item T: 54 months
\end{itemize}

\subsection{Target population}
The target population of this research is women aged between 21-49 years with pain recurrence following previous surgical treatment of endometriosis but wish to avoid removal of ovaries and hysterectomy. The main issue being addressed is endometriosis which is a very common condition in women of reproductive age with about 1 in 10 affected. Endometriosis significantly impacts quality of life because of the side effect it brings, such as pelvic pain. Currently, the standard approach is laparoscopic surgery followed with hormonal contraceptives if there are no plans on conceiving. However, this approach itself doesn't completely remove the chance of pain recurrence in the future and pain recurrence is fairly common with almost half the population being affected. Furthermore, the Scottish national data states that 45\% of women diagnosed with endometriosis has had three or more surgeries. The target population was chosen because there is an increase in awareness on how the current standard laparoscopic surgery approach may not be as clinically and economically effective.

\subsection{Eligibility of the numbers in the trial}
The trial design plans to randomise 400 women to 2 groups: intervention (use of GnRH analogues with add-back HRT) and comparison (laparoscopic surgery) but the primary outcome data needs data on 320 women two years post randomisation. This is due to potential attrition that was set at 20\% which can be verified by the following
\begin{equation*}
    400 - (400 \times 20\%) = 400 - 80 = 320
\end{equation*}
The numbers required are consistent.

\subsection{Informed consent and baseline measures}

\end{document}